%! TEX root = preface.tex
% the main TeX file which is intended to compile, :VimtexReload after adjustment
   % this file should be included in the main file
% :h vimtex-tex-root
\documentclass[../riset_senja.tex]{subfiles}
%ensure the class of docs
\begin{document}

% kebutuhn waterfont sbg multiguna
Kawasan tepi laut kini mengalami perkembangan yang pesat. Banyak tepi laut yang menunjukkan perubahan yang signifikan. Orang-orang dengan latar belakang yang berbeda mulai menyadari keberadaan ruang yang satu ini. Sedikit atau banyak telah memberikan manfaat bagi penggunanya.

% wf strategis
Peningkatan popularitas kawasan tepi laut telah menciptakan  keragaman istilah. Masyarakat lokal sering menggunakan istilah tepi laut, pantai, pesisir dan pinggir laut, serta istilah umum lakefront dan riverfront. Namun pada penelitian ini, istilah tepi laut merujuk pada area pesisir. Area yang berpotensi menjadi pusat aktivitas ekonomi dan masyarakat yang strategis melalui sumber daya, lingkungan binaan, warisan budaya, serta kondisi iklim yang unggul \citep{bhagwat1975,hoppe2011}.

Sebagian besar masyarakat menempatkan ruang publik sebagai wahana rekreasi. Disini muncul jenis tepi laut seperti

Terlepas dari semakin pentingnya ruang tepi danau dalam perencanaan kota modern, ketidaksetaraan spasial dan sosial yang terus-menerus dalam aksesibilitas dan penggunaannya menyoroti tantangan signifikan dalam mencapai distribusi yang merata di berbagai kelompok demografis.

Pemanfaatan ruang tepi air yang optimal tidak semata ditentukan oleh kualitas fisik kawasan, melainkan oleh sejauh mana infrastruktur sosial dan pengembangan guna lahan campuran diintegrasikan untuk memperluas keterjangkauan, meningkatkan intensitas kunjungan, dan pada akhirnya menopang kesejahteraan masyarakat secara berkelanjutan


% Sejumlah daerah bahkan mengalami peningkatan jumlah ruang publik berjenis tepi laut. Beberapa mengalami perubahan secara signifikan, mencerminkan pergeseran paradigma dari sekadar kawasan industri menjadi ruang multifungsi yang lebih berorientasi pada manusia.

Ruang publik tepi laut merupakan elemen yang sangat adaptif terhadap kebutuhan penggunanya. Fleksibilitas ini menciptakan tata ruang yang dinamis dan membentuk kualitas ruang yang beragam.
% klasifikasi xing wd,wr,dan we
\citet{xing2025} menemukan tiga macam tipologi pemanfaatan tepi laut berdasarkan ketergantungannya terhadap air (waterfront dependence), yakni tepi laut bergantung pada air (water-dependent), tepi laut terkait air (water-related), dan tepi laut ditingkatkan oleh air (water-enhanced).
% industri maritim dan publik space
Di sisi lain, \citet{hoppe2011} mengelompokkan kegiatan tepi laut ke dalam kategori industri maritim \textit{water-dependent} dan ruang publik (water-enhanced).
Selain itu, \citet{shi2025} mengklasifikasikan tipologi ruang tepi danau perkotaan (urban lakefront spaces) berdasarkan pola temporal dan intensitas kunjungan (exposure) ke dalam empat kategori: urban lifelines yang ramai baik pada hari kerja maupun akhir pekan, weekend havens yang hanya ramai di akhir pekan, workday blues yang didominasi kunjungan pada hari kerja, serta quiet waters yang memiliki tingkat kunjungan rendah sepanjang minggu.




Ruang publik tepi laut merupakan elemen yang sangat adaptif terhadap kebutuhan penggunanya. Fleksibilitas ini menciptakan tata ruang yang dinamis dan membentuk kualitas ruang yang beragam.


Karakteristik ini tidak hanya ditentukan oleh elemen fisik, tetapi juga oleh interaksi sosial yang terjadi di dalamnya, menjadikannya entitas yang terus berkembang seiring dengan perubahan pola perilaku masyarakat.





Ruang publik tepi laut yang selanjutnya disebut tepi laut adalah sebuah penyedia yang sangat lentur.  Mereka menanggapi masukan dari para pengguna.

Sehingga memunculkan keanekaragaman ruang yang mampu meningkatkan kualitas hidup suatu wilayah pada sebuah tepi laut.
\citet{xing2025} menemukan bahwa bentuk tepi laut



Mereka merespon dari penggunanya. Tidak sedikit ruang memiliki kualitas ganda saat ini.

Mereka merespon dari penggunanya.




Pemanfaatan potensi lanskap pesisir ini tidak hanya berfokus pada estetika visual, tetapi juga harus mengintegrasikan fungsi ekologis dan sosial agar tercipta ruang publik yang berkelanjutan bagi masyarakat setempat.






% pemanfaatan
Pemanfaatan ini memicu konflik lahan, terutama antara kebutuhan pengembangan cara hidup konvensional dan tekanan urbanisasi untuk penyediaan ruang publik. Pengembangan kawasan pesisir secara konvensional mencakup sektor perumahan, perikanan tangkap, dan industri maritim. \citet{rachmadyanti2026} mengklashijau yang direstorasi secara aktif berpartisipasi dalam pembentukan kotaifikasikan pemanfaatan lahan di tepi laut ke dalam tiga kategori: perumahan, kawasan campuran komersial-perumahan, serta aktivitas komersial pariwisata.



Dinamika ini menuntut pendekatan perencanaan yang lebih komprehensif, di mana identitas kawasan tidak hanya dipandang sebagai elemen estetika, melainkan sebagai fondasi utama dalam menciptakan ruang publik yang berkelanjutan dan responsif terhadap kebutuhan sosial masyarakat.





%====================================
% Batas Ediiting
%====================================

Ketegangan muncul ketika alokasi ruang untuk sektor-sektor tersebut bersinggungan dengan kebutuhan akan ruang terbuka hijau dan aksesibilitas publik yang inklusif, yang sering kali terpinggirkan oleh kepentingan ekonomi yang lebih dominan.




Kawasan tepi laut berkembang pesat karena tingginya kebutuhan masyarakat akan ruang multifungsi yang mengintegrasikan rekreasi, belanja, dan budaya, sehingga mampu meningkatkan kualitas hidup wilayah tersebut.





Perkembangan tersebut memunculkan keanekaragaman ruang \citep{huang2025}.




Pesatnya perkembangan kawasan tepi laut akibat tingginya kebutuhan masyarakat akan ruang aktivitas yang beragam menuntut pemahaman mendalam mengenai identitas kawasan tersebut.





Keanekaragaman karakteristik wilayah pesisir menyulitkan penentuan arah pengembangan yang berdampak optimal bagi masyarakat \citep{hoppe2011}.





\citet{hoppe2011} meneliti sebuah kawasan tepi laut yang berfungsi ganda yakni ruang publik dan kawasan industri.












% research gap
The research on public space identity has been conducted numerous times \citep{oktay2015}. However, only a few studies have been carried out in waterfronts to make this place more successful.
This public space is more valuable because it has distinguished characteristics gained from its proximity to the water. Also, its significant contribution towards the city development.


---
Berdasarkan abstrak makalah, "Penggunaan lahan multifungsi dalam pembaruan kawasan pelabuhan: pola sebaran fisik fungsi perkotaan" oleh Antoni Remesar dan Jo PedroO Pedro Costa, berikut adalah temuan dan konsep utama:

Pembaruan Pelabuhan sebagai Lab Perkotaan: Proyek pembaruan perkotaan di port Adelaide bertindak sebagai laboratorium dunia nyata untuk desain perkotaan kontemporer, berfungsi sebagai lokasi strategis yang menerima dana signifikan dan menampilkan pola penggunaan lahan multi-keluarga.

Tipologi Fungsi Perkotaan: Penulis mengkategorikan fungsi dalam proyek pembaruan pelabuhan menjadi dua tipe struktural utama:

Fungsi Perkotaan yang Dominan: Zona kunci yang menjadi ciri dan mendominasi area tertentu (mis., perumahan,komersial / perkantoran, industri, utilitas publik, atau area penggunaan khusus). Zona-zona luas ini dapat menampung fungsi-fungsi lokal yang lebih kecil di dalam batas-batasnya.

Fungsi Perkotaan yang Terletak: Fungsi spesifik dan non-dominan yang terletak di dalam batas area dominan untuk melayani kebutuhan yang ditargetkan (mis., sekolah, museum, toko lokal, dan layanan publik).

Klasifikasi Pola: dengan melakukan studi kasus komparatif di beberapa energi terbarukan tepi laut Eropa (termasuk Amsterdam, Rotterdam, Oslo, dan Barcelona), penelitian ini menunjukkan bahwa distribusi fisik dan kombinasi dari fungsi dominan dan lokasi ini mengikuti pola desain perkotaan yang berbeda dan dapat digeneralisasikan daripada terjadi secara acak.


Harbour Renewal as Urban Labs: Urban renewal projects in port areas act as real-world laboratories for contemporary urban design, serving as strategic sites that receive significant financing and display patterns of multifunctional land use.

Typology of Urban Functions: The authors categorize the functions within harbour renewal projects into two main structural types:

Dominant Urban Functions: Key zones that characterize and dominate a given area (e.g., residential, commercial/office, industrial, public equipment, or special use areas). These broad zones may house smaller, localized functions within their boundaries.

Located Urban Functions: Specific, non-dominant functions situated inside the perimeter of a dominant area to serve targeted needs (e.g., schools, museums, local shops, and public services).

Pattern Classification: By conducting comparative case studies across several European waterfront renewals (including Amsterdam, Rotterdam, Oslo, and Barcelona), the study demonstrates that the physical distribution and combination of these dominant and located functions follow distinct, generalizable urban design patterns rather than occurring at random.

---
Argumen Kunci
Tantangan Tepi Air dengan Kepadatan Sangat Tinggi: Para penulis berpendapat bahwa tepi air bersejarah dengan kepadatan tinggi menghadapi fragmentasi spasial yang akut, penuaan bangunan, pencemaran lingkungan, dan penurunan fungsional karena ketidaksesuaian antara penggunaan lahan pelabuhan tradisional dan perkembangan pariwisata/permainan modern.

Di luar batas urbanisasi Pasca-Industri / Vertikal: Sementara pembaruan tepi laut global sering kali mengikuti transformasi komersial / rekreasi atau pembangunan vertikal bertingkat tinggi, pendekatan semacam itu di zona bersejarah dengan kepadatan tinggi berisiko terhadap gentrifikasi ekologis, ruang publik yang diprivatisasi, dan ketidaksetaraan sosial jika warisan budaya lokal dan kebutuhan masyarakat diabaikan.

Perlunya Kerangka Analisis Terintegrasi: Makalah ini berpendapat bahwa hanya mengandalkan metrik spasial fisik atau data sosial ekonomi tidak memadai; menggabungkan data POI (aktivitas fungsional) dengan Sintaksis spasial (konfigurasi spasial) memberikan diagnosis yang lebih tepat tentang defisiensi fungsional dan ketidaksesuaian spasial.

Pembaruan Seimbang \& Keadilan Spasial: para penulis memperdebatkan strategi pembaruan yang berakar pada keadilan spasial dan ketahanan ekologis-mengadvokasi peningkatan infrastruktur jalan kaki publik, pelestarian warisan pelabuhan, dan integrasi budaya Portugis/lokal yang lebih baik ke dalam praktik komersial daripada komersialisasi murni.

Challenges of Ultra-High-Density Waterfronts: The authors argue that high-density historical waterfronts face acute spatial fragmentation, building aging, environmental pollution, and functional decline due to a mismatch between traditional port land use and modern tourism/gaming developments.

Beyond Post-Industrial / Vertical Urbanism Limitations: While global waterfront renewal often follows commercial/recreational transformation or high-rise vertical developments, such approaches in high-density historic zones risk ecological gentrification, privatized public space, and social inequality if local cultural heritage and community needs are ignored.

Necessity of an Integrated Analytical Framework: The paper contends that relying solely on physical spatial metrics or socio-economic data is insufficient; combining POI data (functional activity) with Space Syntax (spatial configuration) provides a more precise diagnosis of functional deficiencies and spatial mismatches.

Balanced Renewal \& Spatial Justice: The authors argue for renewal strategies rooted in spatial justice and ecological resilience—advocating for enhanced public walking infrastructure, preservation of port heritage, and better integration of Portuguese/local culture into commercial practices rather than pure commercialization.



Ruang arsitektural mengintegrasikan berbagai fungsi, termasuk ruang makan, perdagangan, hiburan, dan kantor [29]. Pembangunan kembali ini telah mengubah kawasan tepi laut menjadi pusat pertumbuhan ekonomi yang vital di jantung pembangunan perkotaan. Distrik dermaga Semenanjung Makau menonjol karena warisan pelabuhannya yang kaya dan akar budaya Portugis. Namun, perkembangan komersial saat ini gagal menyelaraskan secara efektif dengan sumber daya budaya dan pariwisata, yang mengakibatkan ketidaksesuaian antara jenis bisnis dan kebutuhan spasial. Ke depan, Makau harus memanfaatkan pengalaman internasional untuk mengeksplorasi model yang memadukan revitalisasi warisan dengan penawaran komersial yang khas, seperti masakan Portugis dan pasar artisanal.

The architectural spaces integrate multiple functions, including dining, commerce, entertainment, and office spaces [29]. This redevelopment has transformed waterfront areas into vital economic growth hubs at the heart of urban development. The Macau Peninsula pier district stands out for its rich port heritage and Portuguese cultural roots. However, current commercial developments have failed to effectively align with cultural and tourism resources, resulting in a mismatch between business types and spatial needs. Moving forward, Macau should draw on international experience to explore models that blend heritage revitalization with distinctive commercial offerings, such as Portuguese cuisine and artisanal markets.



% Keberadaan identitas yang kuat pada suatu kawasan tidak hanya mendukung terwujudnya keragaman, keberlanjutan, dan inklusivitas, tetapi juga menjadi motivasi bagi masyarakat untuk mengunjungi kawasan tepi laut \citep{petrtylova2022}.
Ini mendukung peningkatan populasi yang terjadi di perkotaan.




Keberadaan identitas yang kuat pada suatu kawasan tidak hanya mendukung terwujudnya keragaman, keberlanjutan, dan inklusivitas, tetapi juga menjadi motivasi bagi masyarakat untuk mengunjungi kawasan tepi laut \citep{petrtylova2022}.

Keberadaan identitas yang kuat pada suatu kawasan tidak hanya mendukung terwujudnya keragaman, keberlanjutan, dan inklusivitas, tetapi juga menjadi motivasi bagi masyarakat untuk mengunjungi kawasan tepi laut.

Pendekatan sensitif terhadap pengembangan wilayah berkarakter unik di tepi laut dapat menghadirkan keragaman, keberlanjutan, dan inklusivitas, serta memberikan alasan untuk mengunjungi ruang sungai."

A sensitive approach to the development of unique character areas of ▎ waterfronts can bring diversity, sustainability and inclusivity, and provide reasons to visit the river space."


Penelitian ini bertujuan mengungkap identitas yang membentuk suatu ruang.
The purpose of this research is to uncover the identity that defines a space.


\bibliographystyle{apalike}
\bibliography{../senja.bib} % required for citecompletion, not work in mainfile
\end{document}

